\begin{example}\label{rl-0008}
  Let \(\mathfrak M\) denote the stack of nodal curves of some fixed genus (we suppress little \(g\) from the notation since we want to use it to denote a group element) and endow it with the trivial action of an algebraic group \(G\). Consider a scheme \(B\). Specializing Remark \ref{rmk:description-of-homotopy-fixed-stack} to this situation, we see that the points of \(\mathfrak M^G(B)\) are given by pairs \((x,\alpha)\) where \(x\in \mathfrak M(B)\) and \(\alpha:G(B)\to \Aut(x)\) is a group homomorphism. The point \(x\) defines a family of nodal curves \(E\to B\) over \(B\) by pulling back the universal curve over \(\mathfrak M\), and \(\alpha\) gives an automorphism
  \begin{center}
    \begin{tikzcd}
      E\arrow[r, "\sim"] \arrow[rd] & E \arrow[d]\\
      & B
    \end{tikzcd}
  \end{center}
  of this family for every element of \(G(B)\). Thus, \(\mathfrak M^G\) parameterizes families of nodal curves with an action of \(G_B\) on \(E\) over \(B\), and \(\eta:\mathfrak M^G\to \mathfrak M\) is the map that forgets this \(G\)-action. Whenever there is such a family \(E\to B\) with a nontrivial \(G_B\)-action on \(E\), \(\eta\) will not be a monomorphism.
\end{example}
